%%%%%%%%%%%%%%%%%%%%%%%%%%%%%%%%%%%%%%%%%%%%%%%%%%%%%%%%%%%%%%%%%%%%%%%%%%%
% Función: guía de los formularios de TFG implementados para la EPS-UAH.
%
% Uso: consulta el PDF para preparar los formularios. Normalmente no debes
% modificar este fichero, salvo para mantener la documentación de la plantilla.
%%%%%%%%%%%%%%%%%%%%%%%%%%%%%%%%%%%%%%%%%%%%%%%%%%%%%%%%%%%%%%%%%%%%%%%%%%%

\documentclass[10pt,a4paper]{article}
\usepackage[T1]{fontenc}
\usepackage[utf8]{inputenc}
\usepackage{lmodern}
\usepackage[spanish,es-lcroman]{babel}
\usepackage[margin=2cm]{geometry}
\usepackage{enumitem,tabularx,booktabs}
\usepackage[colorlinks=true,urlcolor=blue,linkcolor=blue]{hyperref}
\usepackage{microtype}
\setlength{\parindent}{0pt}
\setlength{\parskip}{4pt}
\raggedright
\setlength{\emergencystretch}{2em}
\setlist[itemize]{leftmargin=*,nosep,topsep=3pt}
\newcommand{\variable}[1]{\texttt{\textbackslash{}#1}}
\newcommand{\fichero}[1]{\path{#1}}
\title{Guía del papeleo de TFG\\\large Escuela Politécnica Superior --- Universidad de Alcalá}
\author{}
\date{}

\begin{document}
\begin{center}
{\LARGE Guía del papeleo de TFG\par}
\medskip
{\large Escuela Politécnica Superior --- Universidad de Alcalá\par}
\end{center}

\section{Qué encontrarás aquí}

Hemos preparado ocho formularios que reutilizan los datos de tu trabajo para que no tengas que introducirlos una y otra vez. En esta guía te contamos para qué sirve cada uno, qué variables utiliza y qué tendrás que revisar antes de presentarlo. Se refiere a la normativa de TFG de la EPS-UAH, aprobada el 25/06/2024 y revisada el 05/06/2026; no pretende cubrir todos los centros de la UAH ni sustituir los formularios y procedimientos oficiales vigentes.

Puedes localizar los documentos mediante este resumen. Cada nombre enlaza con su explicación dentro de la guía.

\begin{tabularx}{\textwidth}{@{}Xl@{}}
\toprule
Documento & Anexo EPS-UAH \\
\midrule
\hyperref[form:conjunta]{Autorización conjunta de publicación} & VIII \\
\hyperref[form:autor]{Autorización de la persona autora} & VII \\
\hyperref[form:extranjero]{Autorización del tutor extranjero} & IX \\
\hyperref[form:confidencialidad]{Compromisos de confidencialidad} & VI \\
\hyperref[form:convenio]{Convenio de cooperación educativa} & V \\
\hyperref[form:informe]{Informe del tutor y rúbrica} & III \\
\hyperref[form:plagio]{Informe de control de plagio} & III \\
\hyperref[form:acta]{Acta de defensa y rúbrica del tribunal} & X \\
\bottomrule
\end{tabularx}

\section{Antes de generar los documentos}

Abre \fichero{Config/myconfig.tex} y sustituye los datos de ejemplo por los tuyos. Selecciona una titulación de grado de la EPS-UAH mediante \variable{myDegree} y el idioma mediante \variable{myLanguage}. El identificador del grado determina la universidad, el centro y el tipo de trabajo; no debes editar las variables derivadas \variable{myDegreefull}, \variable{mySchool}, \variable{myUniversity} o \variable{myWorkTypeFull} para rellenar un formulario.

Tu nombre completo, \variable{myAuthorFullName}, se obtiene normalmente de \variable{myAuthorName} y \variable{myAuthorSurname}. Cuando una ficha mencione el nombre completo, revisa estas dos variables. Define \variable{myBookTitleSpanish} y, cuando proceda, \variable{myBookTitleEnglish}: \variable{myBookTitle} selecciona el título según el idioma. Las fichas siguientes indican los datos particulares de cada formulario, incluidos los utilizados por sus fragmentos auxiliares.

La fecha \variable{myPaperworkDate} toma por defecto el valor de \variable{myThesisDepositDate}; la versión inglesa hace lo mismo con \variable{myThesisDepositDateEnglish}. Puedes modificar esas fechas sin cambiar los comandos del formulario. No confundas estas variables con la fecha de firma digital o con la fecha del equipo, utilizadas en algunos documentos.

\subsection{Género gramatical y singular o plural}

La plantilla adapta expresiones como \emph{autor/autora}, \emph{tutor/tutora}, \emph{el/la} o \emph{D./Dª.} mediante las variables de género. Debes escribir exactamente \texttt{male} o \texttt{female}; el nombre de la persona no se utiliza para deducirlo y la implementación no valida otros valores. Comprueba especialmente \variable{myAuthorGender}, \variable{myAcademicTutorGender} y \variable{myCoTutorGender}. Las fichas te indican cuáles afectan realmente a cada documento.

Si no tienes cotutor, deja \variable{myCoTutorFullName} vacío: \verb|\newcommand{\myCoTutorFullName}{}|. Su presencia controla las secciones y firmas opcionales y, donde se usan los comandos correspondientes, el singular o plural. No escribas \texttt{N/A} como nombre para indicar su ausencia, porque entonces la plantilla lo tratará como una persona configurada.

Por ejemplo, la autorización conjunta puede mostrar \emph{la autora} y \emph{la tutora} si ambas variables de género son \texttt{female}. En los compromisos individuales se obtiene \emph{El Profesor} o \emph{La Profesora}; el tutor \emph{tutoriza}, el cotutor \emph{cotutoriza} y un miembro del tribunal \emph{evaluará}. No todos los textos se adaptan: algunos rótulos institucionales son fijos. Los comandos comunes de plural para dos supervisores utilizan actualmente \emph{tutores/directores}, también si ambas personas son mujeres.

\subsection{Casillas configurables y casillas fijas}

No todas las casillas se controlan desde \fichero{Config/myconfig.tex}. Las fichas siguientes distinguen las que dependen de una variable de las que están escritas directamente en el formulario. Respeta las mayúsculas de los valores: \texttt{Yes}/\texttt{No} para publicación y \texttt{YES}/\texttt{NO} para la propuesta de Matrícula de Honor. Tras cambiar la configuración, vuelve a compilar y comprueba las marcas en el PDF.

Si una casilla es fija, debes editar el fichero \texttt{.tex} del formulario cuando su selección no corresponda a tu caso: \verb|$\XBox$| imprime una casilla marcada y \verb|$\Box$| una vacía. Intercambia ambos comandos para cambiar la respuesta y deja marcada solo la opción que proceda. Cambiar una variable de publicación no modifica las casillas fijas.

\section{Publicación y confidencialidad}

\subsection{Autorización conjunta de publicación}\label{form:conjunta}

\fichero{TFG-AutorizacionALLPubAbierto-EPS-UAH.tex}. Anexo VIII de la normativa EPS-UAH, revisada el 05/06/2026. Reúne la autorización y las firmas del estudiante, el tutor y, si existe, el cotutor. No es necesario generar también la autorización individual del autor salvo que el procedimiento aplicable te la pida.

Configura los siguientes datos:
\begin{itemize}
  \item Estudiante y título: \variable{myAuthorFullName}, \variable{myAuthorDNI}, \variable{myBookTitleSpanish} y \variable{myDegree} (para la titulación).
  \item Tutor: \variable{myAcademicTutorFullName}, \variable{myAcademicTutorDNI} y \variable{myAcademicTutorDepartmentOrInstitution}.
  \item Cotutor, si existe: \variable{myCoTutorFullName}, \variable{myCoTutorDNI} y \variable{myCoTutorDepartmentOrInstitution}.
  \item Publicación: \variable{myAuthorizationOpenPublishing}, con \texttt{Yes} o \texttt{No}. Este formulario no muestra opciones de embargo.
  \item Género: \variable{myAuthorGender}, \variable{myAcademicTutorGender} y \variable{myCoTutorGender}; adaptan los rótulos y las firmas. La firma del cotutor desaparece si su nombre está vacío.
\end{itemize}
Las casillas dependen de \variable{myAuthorizationOpenPublishing}: \texttt{Yes} marca \emph{AUTORIZAN} y deja vacía \emph{NO AUTORIZAN}; \texttt{No} hace lo contrario. \variable{myAuthorizationOpenPublishingEmbargoMonths} no interviene en este formulario. Revisa la decisión de autorización y recoge las firmas correspondientes.

\subsection{Autorización de la persona autora}\label{form:autor}

\fichero{TFG-AutorizacionAutorPubAbierto-EPS-UAH.tex}. Anexo VII de la normativa EPS-UAH, revisada el 05/06/2026. Contiene la declaración y autorización individual de quien ha realizado el trabajo y el espacio de firma por la Universidad; no utiliza los datos del tutor ni del tribunal.

Los datos utilizados son \variable{myAuthorFullName}, \variable{myAuthorDNI}, \variable{myBookTitle} (según \variable{myLanguage}), \variable{myPaperworkDate}, \variable{myResearchVicerrector} y la universidad derivada de \variable{myDegree}. Revisa \variable{myAuthorGender} y \variable{myResearchVicerrectorGender}: aparecen expresiones como \emph{autor/autora}, \emph{notificado/notificada} y \emph{Vicerrector/Vicerrectora}. La declaración no utiliza las casillas de \variable{myAuthorizationOpenPublishing}: revisa el texto y las firmas antes de utilizarla.

\subsection{Autorización del tutor extranjero}\label{form:extranjero}

\fichero{TFG-AutorizacionTutorExtranjeroPubAbierto-EPS-UAH.tex}. Anexo IX de la normativa EPS-UAH, revisada el 05/06/2026. Está redactado en inglés y utiliza los datos del \emph{tutor académico}, no los del cotutor: revisa que esta correspondencia sea adecuada para tu caso.

Configura \variable{myAcademicTutorFullName}, \variable{myAcademicTutorForeignUniversity}, \variable{myAcademicTutorForeignCity}, \variable{myAuthorFullName}, \variable{myBookTitleEnglish} y \variable{myPaperworkDateEnglish}. El texto utiliza \emph{Professor} y no emplea variables de género. Las casillas son fijas: \emph{I agree} está marcada y \emph{I disagree} está vacía; no dependen de \variable{myAuthorizationOpenPublishing}. Si el tutor no autoriza, debes editar este fichero e intercambiar \verb|$\XBox$| y \verb|$\Box$| en esas dos opciones. Comprueba la decisión y obtén la firma y el sello indicados.

\subsection{Compromisos de confidencialidad}\label{form:confidencialidad}

\fichero{TFG-CompromisoConfidencialidad-EPS-UAH.tex}. Anexo VI de la normativa EPS-UAH, revisada el 05/06/2026. Genera un compromiso para el tutor y otro para cada persona opcional cuyo nombre esté configurado: cotutor, presidente, dos vocales y suplente. El resultado es un único PDF con hasta seis compromisos; en cada uno firman la persona correspondiente y el estudiante.

Revisa estas variables:
\begin{itemize}
  \item Estudiante y trabajo: \variable{myAuthorFullName}, \variable{myAuthorDNI}, \variable{myAuthorGender}, \variable{myBookTitle} y \variable{myPaperworkDate}.
  \item Tutor: \variable{myAcademicTutorFullName}, \variable{myAcademicTutorDNI} y \variable{myAcademicTutorGender}.
  \item Cotutor: \variable{myCoTutorFullName}, \variable{myCoTutorDNI} y \variable{myCoTutorGender}.
  \item Presidente: \variable{myTribunalPresident}, \variable{myTribunalPresidentDNI} y \variable{myTribunalPresidentGender}.
  \item Vocal primero: \variable{myTribunalFirstSpokesperson}, \variable{myTribunalFirstSpokespersonDNI} y \variable{myTribunalFirstSpokespersonGender}.
  \item Vocal segundo: \variable{myTribunalSecondSpokesperson}, \variable{myTribunalSecondSpokespersonDNI} y \variable{myTribunalSecondSpokespersonGender}.
  \item Suplente: \variable{myTribunalAlternateMember}, \variable{myTribunalAlternateMemberDNI} y \variable{myTribunalAlternateMemberGender}.
\end{itemize}
Un nombre opcional vacío omite esa persona y un DNI de profesor vacío deja una línea para completarlo. La fecha y el título siguen los valores de \variable{myPaperworkDate} y \variable{myBookTitle}; este último depende del idioma. El formulario no se activa mediante \variable{myConfidentialContent}: genéralo únicamente si lo necesitas.

\section{Cooperación educativa}

\subsection{Anexo al convenio}\label{form:convenio}

\fichero{TFG-ConvenioCooperacion-EPS-UAH.tex}. Anexo V de la normativa EPS-UAH, revisada el 05/06/2026. Describe el trabajo que se realizará en una entidad externa. La plantilla representa al tutor de la entidad mediante el cotutor.

Utiliza los siguientes datos:
\begin{itemize}
  \item Estudiante: \variable{myAuthorFullName}, \variable{myAuthorDNI}, \variable{myAuthorTelephone} y \variable{myAuthorEmail}.
  \item Tutor académico: \variable{myAcademicTutorFullName}, \variable{myAcademicTutorDNI}, \variable{myAcademicTutorDepartmentOrInstitution} y \variable{myAcademicTutorEmail}.
  \item Entidad y tutor externo: \variable{myCoTutorFullName}, \variable{myCoTutorDepartmentOrInstitution}, \variable{myCoTutorPosition}, \variable{myCoTutorDNI} y \variable{myCoTutorEmail}. Sin cotutor, los campos derivados de él muestran \texttt{N/A}.
  \item Trabajo: \variable{myDegree}, \variable{myBookTitleSpanish} y \variable{myThesisAcademicYear}. La titulación seleccionada proporciona el centro y el tipo de trabajo.
  \item Género: \variable{myAuthorGender} y \variable{myCoTutorGender}, para el tratamiento del estudiante y el rótulo del tutor externo. Algunos rótulos como \emph{tutor/a} permanecen fijos.
\end{itemize}
Completa los campos todavía marcados como \texttt{N/A}: identificación y dirección de la entidad, teléfonos, fechas de inicio y fin, horario y resumen del trabajo. No tienen variables propias en \fichero{myconfig.tex}; puedes completar el formulario o adaptar sus campos en el fuente. La fecha impresa utiliza \variable{today}, es decir, la fecha del equipo durante la compilación, no \variable{myPaperworkDate}. Revisa también las tres firmas: coordinación del trabajo, tutor de la entidad y estudiante.

\section{Informes y evaluación}

\subsection{Informe del tutor y rúbrica}\label{form:informe}

\fichero{TFG-InformeTutorYRubrica-EPS-UAH.tex}. Anexo III de la normativa EPS-UAH, revisada el 05/06/2026. Incluye el informe de plagio, los datos y firmas del tutor y del cotutor, si existe, y la rúbrica del tutor.

El fragmento de plagio utiliza \variable{myAuthorFullName}, \variable{myAuthorDNI}, \variable{myAuthorEmail}, \variable{myDepartment}, \variable{myDegree}, \variable{myBookTitleSpanish} y, con \variable{myLanguage} igual a \texttt{english}, \variable{myBookTitleEnglish}. El informe añade \variable{myAcademicTutorFullName}, \variable{myAcademicTutorPosition}, \variable{myAcademicTutorEmail}, \variable{myCoTutorFullName}, \variable{myCoTutorPosition}, \variable{myCoTutorEmail}, \variable{myThesisAcademicYear} y \variable{myAcademicTutorGrade} (hasta siete puntos). Los rótulos y firmas dependen de \variable{myAcademicTutorGender} y \variable{myCoTutorGender}.

Completa los resultados del control de plagio, las observaciones, la conclusión, los comentarios del tutor y la rúbrica. La fecha indicada es la de la firma digital; \variable{myThesisDefenseDate} no se imprime actualmente en este documento.

\subsection{Informe de control de plagio}\label{form:plagio}

\fichero{TFG-InformeTutorPlagio-UAH.tex}. Es la parte de informe de plagio del anexo III de la normativa EPS-UAH, revisada el 05/06/2026. Conserva el texto de protección de datos UAH de su fragmento auxiliar; no es un nuevo anexo independiente ni incorpora la rúbrica.

Utiliza \variable{myAuthorFullName}, \variable{myAuthorDNI}, \variable{myAuthorEmail}, \variable{myDepartment}, \variable{myDegree}, \variable{myBookTitleSpanish} y, con \variable{myLanguage} igual a \texttt{english}, \variable{myBookTitleEnglish}. Revisa también \variable{myAcademicTutorFullName}, \variable{myAcademicTutorPosition}, \variable{myAcademicTutorEmail}, \variable{myCoTutorFullName}, \variable{myCoTutorPosition} y \variable{myCoTutorEmail}. Los rótulos del tutor y cotutor se adaptan con \variable{myAcademicTutorGender} y \variable{myCoTutorGender}; no se emplea \variable{myAuthorGender} para los rótulos del estudiante de este informe.

Completa coincidencias, observaciones y conclusión del análisis y recoge las firmas. La fecha es la de la firma digital, no \variable{myThesisDefenseDate}.

\subsection{Acta de defensa y rúbrica del tribunal}\label{form:acta}

\fichero{TFG-ActaYRubricaDefensa-EPS-UAH.tex}. Anexo X de la normativa EPS-UAH, revisada el 05/06/2026. Incluye las calificaciones y observaciones del tribunal y su rúbrica.

Configura \variable{myAuthorFullName}, \variable{myAuthorDNI}, \variable{myDegree}, \variable{myBookTitleSpanish}, \variable{myBookTitleEnglish} (si \variable{myLanguage} es \texttt{english}), \variable{myAcademicTutorFullName}, \variable{myCoTutorFullName}, \variable{myDepartment} y \variable{myThesisAcademicYear}. Para las notas utiliza \variable{myAcademicTutorGrade} (hasta siete puntos), \variable{myTribunalGrade} (hasta tres puntos) y \variable{myTribunalMHProposal} (\texttt{YES} o \texttt{NO}). La plantilla suma las dos notas y obtiene la calificación final y su denominación.

Los rótulos de los tutores se adaptan mediante \variable{myAcademicTutorGender} y \variable{myCoTutorGender}. La rúbrica incluida añade \variable{myThesisDefenseDate}, \variable{myTribunalPresident} y \variable{myTribunalFirstSpokesperson}. Su presidencia utiliza \variable{myTribunalPresidentGender}, mientras que el comando dependiente de \variable{myTribunalFirstSpokespersonGender} mantiene el rótulo \emph{Vocal} para ambos valores. Actualmente esta rúbrica solo incluye esas dos firmas: no imprime el segundo vocal ni el suplente. El acta mantiene el rótulo fijo \emph{Autor}. Completa las observaciones y la rúbrica y comprueba que las firmas corresponden a tu tribunal.

Las casillas de propuesta para Matrícula de Honor dependen de \variable{myTribunalMHProposal}: \texttt{YES} marca \emph{Sí} y deja vacía \emph{No}; \texttt{NO} hace lo contrario. La propuesta no se deduce de las notas ni de su suma: debes configurar la decisión del tribunal y comprobar que cumple los requisitos aplicables.

\section{Cómo generar los documentos}

No necesitas Make. Abre el formulario que necesites en tu editor de LaTeX, mantén la estructura de directorios y compílalo con PDFLaTeX desde \fichero{PapeleoTFG/}. Si tu editor lo requiere, establece esa carpeta como directorio de trabajo. En Overleaf, selecciona el fichero del formulario como documento principal. Estas plantillas no necesitan Biber ni \texttt{makeglossaries}; los fragmentos de rúbrica y secciones compartidas no se compilan por separado.

Si utilizas Make, ejecuta estas órdenes desde \fichero{PapeleoTFG/}; necesitas \texttt{latexmk} además de la distribución de LaTeX:
\begin{minipage}{\linewidth}
\begin{verbatim}
make                         # Todos los formularios
make autorizaciones          # Publicación y confidencialidad
make convenios               # Convenio de cooperación
make evaluacion              # Informes, acta y rúbricas
make TFG-InformeTutorPlagio-UAH.pdf
make guia                    # Esta guía, sin generar formularios
make help                    # Opciones disponibles
\end{verbatim}
\end{minipage}
Cada fichero independiente genera un PDF con su mismo nombre. \texttt{make clean} elimina auxiliares y conserva los PDF; \texttt{make cleanall} elimina también los PDF de formularios y de esta guía. También puedes compilar la guía por separado: no lee tu configuración ni incorpora datos personales.

\section{Antes de presentar los documentos}

Dedica unos minutos a revisar el resultado, aunque se haya generado sin errores:
\begin{itemize}
  \item Comprueba que el formulario y su versión corresponden a tu titulación y al procedimiento que te hayan indicado. No todos los documentos se utilizan en todos los trabajos.
  \item Revisa nombres, DNI, afiliaciones, títulos, fechas, géneros, cotutor y miembros del tribunal. No dejes los datos de ejemplo de la plantilla.
  \item Completa campos vacíos y \texttt{N/A}, análisis de plagio, notas, observaciones y decisiones de publicación. Revisa el género también en los tratamientos escritos dentro de nombres o cargos, que no se reescriben automáticamente.
  \item Obtén las firmas que correspondan y sigue el procedimiento de entrega que te indique la EPS-UAH. La plantilla genera documentos, no los firma ni los presenta por ti.
\end{itemize}
\end{document}
